\chapter{System Architecture and Implementation}
\label{ch:implementation}

\section{A modular architecture behind interfaces}
\label{sec:architecture}

The system is a Python package organised into modules, each behind an abstract interface,
so that a component can be replaced without touching the rest
(Figure~\ref{fig:architecture}). There is an audio source
(\texttt{AudioSource}), a voice activity detector (\texttt{VAD}), a speech recogniser
(\texttt{ASR}), the two end-of-turn branches (\texttt{AcousticEOU} and
\texttt{LexicalEOU}), the fusion (\texttt{Fusion}) and the endpoint controller. This
follows directly from the late-fusion choice of Section~\ref{sec:latefusion}: the
recogniser or the detector can be exchanged, a client to an existing microservice in place
of the in-process model, without the decision logic changing at all.

The decision core is isolated from transport in a per-call session,
\texttt{EouSession}. The session takes raw 16\,kHz mono audio and the latest partial
transcript, and emits a decision at each micro-pause. It is driven by an \texttt{Orchestrator}, which owns an audio source and a recogniser for
the demonstrations and the evaluation harness, and by the WebSocket service, which drives a
session directly.
Both reach the same code, so there is one decision path and not two.

The session carries an explicit statefulness contract, which matters in production:
the heavy models may be shared across concurrent calls, but the VAD holds state and must be
instantiated per session.

\section{The real-time path}
\label{sec:realtime}

The audio source emits frames of 512 samples at 16\,kHz, that is 32\,ms. In production the
transport does not respect that framing, since telephony delivers 20\,ms frames, so the
session accepts a chunk of arbitrary length and re-frames it internally to the fixed size
the VAD expects. Each frame is written to a ring buffer holding the last eight seconds of
signal, passed to the recogniser, and labelled speech or silence by the VAD. When a
speech-to-silence transition lasts beyond 200\,ms, an evaluation is triggered. That is the
micro-pause at which the system decides to end the turn or to wait, and an internal latch
guarantees it fires at most once per turn.

\begin{figure}[H]
\centering
\resizebox{\textwidth}{!}{%
\begin{tikzpicture}[
  font=\small,
  box/.style={draw=black!60, rounded corners=1.5pt, align=center,
              minimum height=9mm, inner xsep=2.5mm, inner ysep=1.5mm},
  fl/.style={-{Stealth[length=2mm]}, draw=black!60},
  sn/.style={-{Stealth[length=2mm]}, draw=black!60, dashed}]

\node[box] (src)  at (0,0)      {Audio source\\\footnotesize 512 samples / 32\,ms};
\node[box] (asr)  at (0,2.0)    {ASR (Azure)\\\footnotesize interim hypotheses};
\node[box] (ring) at (0,-2.0)   {Ring buffer\\\footnotesize last 8\,s};
\node[box] (vad)  at (4.3,0)    {VAD\\\footnotesize speech / silence};
\node[box] (gate) at (8.2,0)    {Micro-pause\\\footnotesize silence $>$ 200\,ms};
\node[box] (ac)   at (12.6,1.6) {Acoustic branch\\\footnotesize Smart Turn v3};
\node[box] (lex)  at (12.6,-1.6){Lexical branch\\\footnotesize CamemBERT};
\node[box] (fus)  at (16.9,0)   {Fusion\\\footnotesize \texttt{RuleFusion}};
\node[box] (ep)   at (20.9,0)   {Endpoint\\\footnotesize controller};
\node      (out)  at (24.2,0)   {\textbf{END} / \textbf{WAIT}};

\draw[fl] (src) -- (vad);
\draw[fl] (src) -- (asr);
\draw[fl] (src) -- (ring);
\draw[fl] (vad) -- (gate);
\draw[fl] (gate) |- (ac);
\draw[fl] (gate) |- (lex);
\draw[fl] (ring) -| (ac);
\draw[sn] (asr)  -| (lex) node[midway, above, font=\footnotesize] {snapshot, never awaited};
\draw[fl] (ac) -| (fus);
\draw[fl] (lex) -| (fus);
\draw[fl] (fus) -- (ep);
\draw[fl] (ep) -- (out);

\begin{scope}[on background layer]
\node[draw=black!35, dashed, rounded corners=2pt, fit=(ac)(lex),
      inner sep=3.5mm, label={[font=\footnotesize, black!55]below:%
      concurrent, in an executor}] {};
\end{scope}
\end{tikzpicture}}
\caption[Pipeline architecture]{The decision pipeline. Every frame feeds the ring buffer,
the recogniser and the VAD. A speech-to-silence transition beyond 200\,ms triggers an
evaluation, in which the two branches run concurrently off the event loop. The lexical branch reads the latest partial transcript as a snapshot, so the decision never waits on
the recogniser; both model calls are awaited, and the added latency is the slower of the
two. The fusion shown is \texttt{RuleFusion}, the variant the system runs; the learned
form of Equation~\ref{eq:fusion} substitutes behind the same interface but is evaluated
offline only.}
\label{fig:architecture}
\end{figure}

The session keeps two clocks. Silence duration is
measured on an audio-sample clock, advanced by the frame duration, so a given recording
always produces the same timing regardless of how fast the machine runs. Wall-clock time is
used for one thing only: measuring how long the models took. A system timing silence by wall clock would behave differently on a loaded machine and on an
idle one, and a replayed clip would not reproduce its own earlier result. This does not
affect Chapter~\ref{ch:results}, whose figures come from stored predictions; it matters for
the live pipeline and the clip-replay harness.

At the moment of evaluation the acoustic and lexical branches are computed concurrently,
off the event loop in an executor, so that incoming frames continue to be processed. The partial transcript is read as a local snapshot and never awaited. That is the
non-blocking property of Section~\ref{sec:nonblocking} in concrete terms: the recogniser is
off the critical path. The two model calls themselves are both awaited, so the added latency
is $\max(\text{acoustic}, \text{lexical})$ and not their sum, and no timeout guards the
lexical future.

\section{The two branches and the recogniser}
\label{sec:branches}

The acoustic branch applies Smart Turn v3 (version 3.2, CPU ONNX export) to the last eight
seconds of the ring buffer and returns a probability that the turn has ended. The lexical
branch applies CamemBERT to the latest partial and returns a completeness probability. The
two variants of Section~\ref{sec:frozen} sit behind the same interface, frozen embeddings
with a classifier for the benchmark and the fine-tuned model for the domain and the live
demonstration. A French heuristic guard (\texttt{FrenchSemanticEOU}) complements the model
on the hard cases, detecting a spelling sequence still in progress, digits, and function
words in final position.

Transcription reuses the production Azure service, with the same language and audio format.
The one non-invasive addition is that it subscribes to the interim hypotheses as well as
the final results, and those feed the lexical branch. Production itself is unchanged.

\section{Fusion and the endpoint controller}
\label{sec:fusionimpl}

\texttt{RuleFusion}, one of the two implementations behind the \texttt{Fusion} interface,
applies a bilateral veto: if either signal falls below 0.25, that is, if either branch clearly reads
mid-turn, the system waits. If both exceed 0.85 it endpoints immediately. Otherwise it
averages the two and hands the result on. \texttt{LogisticFusion} substitutes for it without the
calling code changing. It implements Equation~\ref{eq:fusion} with one addition: a fifth
term on the elapsed silence, normalised by the maximum-silence limit. The benchmark
experiments of Chapter~\ref{ch:results} use the four-parameter form of
Equation~\ref{eq:fusion} without that term, so that the reported gain is attributable to
the two probability signals alone and not to timing information the acoustic baseline never
receives.

One boundary should be drawn here. The learned fusion and the cost rule were staged from
the outset, with the rule-based fusion as the committed deliverable, and neither reached the
running path: \texttt{LogisticFusion} is instantiated nowhere outside a unit test, and no
cost ratio appears anywhere in the runtime package. 

The data strategy of Section~\ref{sec:datastrategy} then fitted both the fusion and the calibrator on generic
French, and Section~\ref{sec:fusionresults} measures what a component fitted on the wrong
domain costs. What the running system demonstrates is the pipeline, its real-time and
non-blocking behaviour, and the capture that produced the in-domain set; the decision layer
of Chapter~\ref{ch:results} is established offline on stored predictions, and carrying it
across is deliberately not claimed here.

The endpoint controller then converts confidence into a duration. Rather than a single
threshold it maps confidence onto how much trailing silence is demanded: 200\,ms above
0.85, 600\,ms above 0.5, 1200\,ms below. Confidence therefore buys speed, and a hesitant
score is made to wait longer rather than being refused outright.

Above all of that sits a hard limit at 2000\,ms, which deliberately overrides the veto. A
caller who genuinely stopped mid-sentence, and whose text will never look complete, would
otherwise be held indefinitely by the lexical branch. This is the one place where the design accepts a wrong cut over an unbounded wait. The
reasoning is about failure modes: an interruption is recoverable in conversation, silence
from the agent is not.

\section{Demonstrations and evaluation tooling}
\label{sec:tooling}

Two demonstrations drive the same core. One runs at the command line, from the microphone or
a wav file replayed as a stream; the other runs in the browser. Both print, at each
decision, the acoustic, lexical and fused probabilities, the decision taken, the observed
silence and the compute latency.

Evaluation rests on two tools: a local harness that replays labelled clips and computes
accuracy, false positives, false negatives and latency, and an adapter that plugs the fusion
into the eot-bench protocol so it can be compared with other systems on identical metrics.

The in-domain scoring path did not survive that route, and the reason cost me a day: the
harness reads audio through the Hugging Face \texttt{Audio} feature, which now requires
\texttt{torchcodec}, which does not load against \texttt{ffmpeg}~8. Rather than pin the
toolchain backwards, a standalone scorer reads the caller wav files with \texttt{soundfile}
and scores each pause with the same models and fold splitting. It produces the figures of
Section~\ref{sec:gold}.

\section{Production instrumentation}
\label{sec:instrumentation}

To build the in-domain set, the production system was instrumented. Capture is gated behind
an organisation-level feature flag. When active, the system records at each micro-pause the
caller-side text only, the trajectory of the recogniser's interim hypotheses, timestamps, a
reference to the audio, and resume or interruption flags. Labels come from weak supervision
with no manual annotation: a pause followed by the caller resuming is a hold, a pause
followed by the agent replying without a resume is an end of turn.

The resulting set is 709 calls placed on a single production organisation, of which roughly
four in five are genuine inbound calls and the rest calls made by internal testers to the
same line. Of those, 636 contribute at least one scored decision pause.
Names, first names and ages are masked in the retained transcripts, and the data stays
within Tennor's infrastructure. Because caller audio is retained and linked to a call identifier, it remains
personal data and is treated as such throughout. Appendix~\ref{app:ethics} sets out the
governance in full.

\section{Testing and verification}
\label{sec:testing}

The system is covered by 45 automated tests, exercising the VAD, the ring buffer, the fusion
and endpoint logic, the lexical rules, the types and the metrics. Integration tests that
need heavy models, credentials or audio hardware skip themselves when those are absent, so
the suite runs anywhere. Every result reported in Chapter~\ref{ch:results} has a
documented command that reproduces it, listed in \texttt{reproduce.md} in the submitted
code corpus.
